% =====================================================================
% Section III -- Setup and Same-Target Temporal Semantics
% Paper: Component-Dependent Temporal Effects in Action-Chunked Robot Policies
% Target: ~720 words. IEEE two-column, 8 pages including references.
% Frozen constitution: no method claim, no additive attribution, no pure
% observation-delay estimand. Attack B is conceded here, in III-B.
% =====================================================================

\section{Setup and Same-Target Temporal Semantics}
\label{sec:setup}

\subsection{Policies, control, and the action contract}

All temporal experiments use fixed trained policies; the execution
interventions do not modify policy weights. Our primary policies are frozen
Action Chunking Transformers (ACT)~\cite{zhao2023act} from the official
LeRobot implementation~\cite{cadene2026lerobot}, with chunk size $100$,
temporal ensembling disabled, and no action smoothing. Evaluation is
simulation-only on LIBERO~\cite{liu2023libero} under relative end-effector
(EEF) control at 20\,Hz. A single frozen SmolVLA
pilot~\cite{shukor2025smolvla} provides only heterogeneous cross-policy scope
evidence under the same intervention semantics.

The $7$-dimensional LIBERO action separates into two groups that we treat
as the experimental factor: a $6$-dimensional arm group spanning dimensions
$0$ through $5$ inclusive and a $1$-dimensional gripper group at dimension
$6$. This split is given by the control interface rather than chosen by us,
and it is the only action-group partition we study.

\subsection{Same-target semantics and the estimand}
\label{sec:same-target}

A chunked policy queried at step $q$ maps its observation to a whole
$7$-dimensional predicted chunk $A_q = \pi(o_q)$ whose entry $A_q[k]$ is the
action the policy predicts for physical time $q+k$. We require every executed
action to be a prediction for the time at which it is actually executed:
\begin{equation}
q + k = t .
\label{eq:same-target}
\end{equation}
Writing the source age as $d = t - q$, Eq.~\eqref{eq:same-target} gives
\begin{equation}
k = d .
\label{eq:coupling}
\end{equation}
Source age and chunk offset are therefore not independent factors. This
coupling is a consequence of the same-target contract, not an artifact of
our implementation: any intervention that decouples them, such as executing
$A_{t-d}[0]$ at time $t$, replays an action the policy intended for an
earlier instant and no longer estimates a prediction for $t$.

We state the resulting estimand explicitly. \emph{The natural same-target
temporal relationship in an action-chunked policy jointly determines source
age and prediction offset; we estimate the combined closed-loop effect of
that relationship per action group, and we do not identify a pure
observation-age causal effect.} Section~\ref{sec:limitations} returns to
this point.

Applying Eq.~\eqref{eq:same-target} independently to each group yields
\begin{align}
a_{\mathrm{arm}}(t)  &= P_{\mathrm{arm}} A_{t-d_{\mathrm{arm}}}[d_{\mathrm{arm}}], \\
a_{\mathrm{grip}}(t) &= P_{\mathrm{grip}} A_{t-d_{\mathrm{grip}}}[d_{\mathrm{grip}}].
\end{align}
Here $P_{\mathrm{arm}}$ selects dimensions $0$ through $5$ inclusive from
the whole $7$-dimensional predicted chunk, and $P_{\mathrm{grip}}$ selects
dimension $6$. Both components remain predictions for the same physical time
$t$; only the temporal relationship used to obtain them differs.

\subsection{The factorial conditions}

Crossing $d_{\mathrm{arm}}, d_{\mathrm{grip}} \in \{0, 20\}$ gives four
cells: \textsc{Fresh} $(0,0)$, \textsc{FO20} $(0,20)$, \textsc{Reverse20}
$(20,0)$, and \textsc{FullOld20} $(20,20)$. The value $d=20$ was fixed by a
preregistered condition and was never adjusted after any outcome was
observed.

\subsection{Fixed-horizon chunk execution and its variants}

As our fixed-horizon chunk-execution reference, we include
\textsc{hard-h16}: the policy is queried once every $\HardHorizon$ steps and both groups
read the current committed chunk at offset $t-q$. Its source ages span
0--15 steps, with an observed mean of approximately $\HardMeanSourceAge$ steps in our
rollouts; \textsc{hard-h16} is thus not a fully reactive baseline.

The C2 H16Arm+FreshGrip decomposition also queries the whole policy at every
controller step. Its gripper uses the current offset-$0$ prediction. The arm
uses the scheduled source
$q_{\mathrm{arm}}=\HardHorizon\lfloor t/\HardHorizon\rfloor$ and offset
$k_{\mathrm{arm}}=t-q_{\mathrm{arm}}$. Both groups therefore satisfy the
same-target contract, while only the arm carries the h16 schedule.

\subsection{Blocks, cohorts, and inference}

A paired block is a (task, initial-state) pair evaluated under every
condition with matched seeds, caps, and preprocessing, so that conditions
differ only in temporal relationship. Success is binary per block.

The confirmation cohort comprises LIBERO-Goal tasks $4,6,7,8,9$ and
LIBERO-10 tasks $0,2,4,6,7$ over initial states $0$--$13$, giving
$\ConfirmFreshBlocks$
paired blocks per condition. These tasks were \emph{unseen to executor
development}; they were \emph{not} held out from policy training, and we
make no generalization claim that depends on the stronger reading.

Paired contrasts use the two-sided exact McNemar test. The exact McNemar test
depends only on discordant blocks; we therefore report their counts
explicitly. We accompany each contrast with a paired 95\% interval on the
risk difference and a task-level cluster bootstrap interval, and we report
the leave-one-task-out sign across all ten tasks. Where the two intervals
disagree we report both and draw no conclusion from the narrower one.
